\documentclass[11pt, aspectratio=169]{beamer}
\usepackage[style=authoryear-icomp, backend=biber]{biblatex}
\usetheme{jcar}

%\includeonlyframes{a} % Para compilar solo en la que se está trabajando

%\addbibresource{refs.bib}

\begin{document}
\author[Juan C. Arboleda Rivera]{Juan Camilo Arboleda Rivera}
\title{Probabilidad}
%\subtitle{}
%\titlegraphic{\includegraphics[width=2cm]{udea}}
\institute[U. de A.]{Instituto de Biología\\
  Facultad de Ciencias Exactas y Naturales\\
  Universidad de Antioquia\\
  juan.arboleda2@udea.edu.co}
\date{v: \today}
%\subject{}

\begin{frame}[plain]
  \maketitle
\end{frame}

\begin{frame}
 \frametitle{Contenidos}
 \tableofcontents
\end{frame}

\section{Probabilidad en Ciencia de datos}

\begin{frame}
%[label=a]
{}{}
 \centering
 Supongamos que queremos saber el peso promedio de la\\
 tortuga verde (\textit{Chelonia mydas})\\
 ¿Podemos conocer ese valor?\\
 \vspace{.2em}
 \includegraphics[height=.75\textheight]{tortuga-verde.jpeg}
\end{frame}

\begin{frame}
%[label=a]
{}{}
 \centering \LARGE
 Cuando tomamos datos, tenemos una \textbf{muestra}, no toda la
 \textbf{población}, y en cada muestra los valores pueden ser diferentes
% \includegraphics[height=.85\textheight]{}
\end{frame}

\begin{frame}
%[label=a]
{¿Por qué es importante la probabilidad en ciencia de datos?}{}
 \centering \small
  Tenemos datos de la expresión de un gen en personas sanas y personas
  enfermas. ¿Podemos concluir que en personas enfermas se expresa más el
  gen? \onslide<2>{¿Qué tan probable es observar una diferencia como la
    encontrada entre dos grupos simplemente por azar?}\\

  \vspace{.3cm}

 \includegraphics[height=.7\textheight]{cd177-covid.png}
\end{frame}

\begin{frame}
%[label=a]
{}{}
 \centering \small
 Si tenemos datos de expresión de organismos normales (Wildtype) y datos de
 expresión de organismos mutantes
 ¿Qué tan probable es que se obtengan los siguientes resultados por azar?
 \\
 \vspace{.3cm}

 \includegraphics[height=.75\textheight]{ligand-binding-site.png}
\end{frame}

\begin{frame}
%[label=a]
{¿Por qué es importante la Probabilidad en Ciencia de datos?}{}
 \centering
 \textbf{Existe variabilidad biológica}\\~\\
 Los organismos no son idénticos; existen diferencias genéticas,
 fisiológicas y ambientales. La probabilidad ayuda a describir y comprender
 esa variación.
% \includegraphics[height=.85\textheight]{}
\end{frame}

\begin{frame}
%[label=a]
{}{}
 \centering
 \textbf{Incertidumbre}\\~\\
 Las mediciones nunca son exactas, incluso un mismo valor (e.g. la longitud
 del pico de un colibrí), será diferente si es medido por diferentes
 personas o incluso en mediciones repetidas por la misma persona. La
 probabilidad ayuda a describir y estudiar la incertidumbre en las medidas.\\
 \vspace{.8em}
 \includegraphics[height=.55\textheight]{Incertidumbre-de-Medida.jpg}
\end{frame}

\begin{frame}
%[label=a]
{¿Por qué es importante la Probabilidad en Ciencia de datos?}{}
 \centering
 \textbf{Análisis de datos experimentales}\\~\\
 En biología se trabaja con muestras (por ejemplo, grupos de pacientes o
 poblaciones de organismos), no siempre con todos los individuos. La
 probabilidad permite estimar si los resultados observados podrían deberse
 al azar o si muestran patrones reales.
% \includegraphics[height=.85\textheight]{}
\end{frame}

\begin{frame}
%[label=a]
{}{}
 \centering
 \textbf{Predicción de fenómenos biológicos}\\~\\
 Los modelos probabilísticos ayudan a predecir eventos como la propagación
 de enfermedades, la respuesta de un paciente a un tratamiento o la
 evolución de una población.
% \includegraphics[height=.85\textheight]{}
\end{frame}

\begin{frame}
%[label=a]
{}{}
 \centering
 \textbf{Aprendizaje automático e inteligencia artificial}\\~\\
 Muchos algoritmos usados en bioinformática, como los modelos de
 clasificación de genes, análisis de imágenes médicas o predicción de
 estructuras de proteínas, utilizan conceptos probabilísticos para aprender
 de los datos.
% \includegraphics[height=.85\textheight]{}
\end{frame}

\begin{frame}
%[label=a]
{}{}
 \centering
 \textbf{Evaluación de riesgos y toma de decisiones}\\~\\
 Permite calcular probabilidades de aparición de enfermedades, identificar
 factores de riesgo y comparar diferentes hipótesis científicas.
% \includegraphics[height=.85\textheight]{}
\end{frame}

\begin{frame}
%[label=a]
{}{}
 \centering
 \textbf{Interpretación de información genética}\\~\\
 En áreas como la genómica, la probabilidad se usa para identificar
 mutaciones relevantes, analizar relaciones evolutivas y estimar la función
 de genes.\\
 \includegraphics[height=.75\textheight]{probabilidad-genética.png}
\end{frame}

\begin{frame}
%[label=a]
{}{}
 \centering
 La probabilidad se usa habitualmente en todas las áreas de la Biología,\\
 por ejemplo en la construcción de árboles filogenéticos\\
 \vspace{.2em}
 \includegraphics[height=.85\textheight]{bayesian_trees.png}
\end{frame}

\begin{frame}
%[label=a]
{}{}
 \centering
 De manera general la \textbf{Probabilidad} nos permite cuantificar qué tan
 plausibles son determinados resultados, lo cual es esencial para comprender
 la Naturaleza.
% \includegraphics[height=.85\textheight]{}
\end{frame}

\section{Introducción a la Probabilidad}

\begin{frame}
%[label=a]
{¿Qué es la Probabilidad?}{}
 \centering
 \begin{block}<1->{Interpretación frecuentista}
 Define la probabilidad como la frecuencia relativa de un evento tras
 repetir un experimento un número grande o infinito de veces bajo las mismas
 condiciones.\\~\\

 \[ P(A) = \dfrac{\text{Número de resultados que hacen que el evento } A
 \text{ suceda}}{\text{Número total de resultados}} \]
 \end{block}

 \pause

 \vspace{.4cm}
 \textbf{Ejemplo:}\\
 La probabilidad de lanzar un dado y que salga el número 5 es 1/6.
\end{frame}

\begin{frame}
%[label=a]
{¿Qué es la Probabilidad?}{}
 \centering

 \begin{block}{Interpretación bayesiana}
 Define la probabilidad como el grado de creencia o certeza sobre un evento
 basándose en el conocimiento previo y la nueva evidencia
 \end{block}

 \vspace{.4cm}

 \pause
 \textbf{Ejemplo:}\\
 La probabilidad de que llueva en la noche es del 20\%.
% \includegraphics[height=.85\textheight]{}
\end{frame}

\begin{frame}
%[label=a]
{}{}
 \centering \large
 La función de probabilidad $ P $, asigna a cada evento una probabilidad de
 ocurrencia. $ P(A) $, es la probabilidad de que el evento $ A $ ocurra al
 efectuar un experimento\\~\\

 Por convención la probabilidad es un número entre 0 y 1 o entre 0 y 100\%
 \\~\\

 Un evento imposible tiene probabilidad cero\\
 Un evento que ocurre con certeza tiene probabilidad 1\\
 \vspace{.4em}
	\includegraphics[width=.6\linewidth]{probability_line.png}
\end{frame}

\begin{frame}
%[label=a]
{Espacio muestral}{}
 \centering \Large
 Para el calculo de probabilidades, se suele definir primero el
 espacio muestral ($ E $), definido como {\color{blue} el conjunto de todos
 los resultados posibles de un experimento}
\end{frame}

\begin{frame}
%[label=a]
{Ejemplo}{Lanzamiento de una moneda}
 \centering \Large
 \includegraphics[height=.5\textheight]{moneda500.jpg}\\
 Para una moneda que tiene cara ($ C $) y sello ($ S $), el espacio muestral
 sería: $E=\{C, S\}$.\\

 Además

 $P(E) = P(C) + P(S)  = 1 $
\end{frame}

\begin{frame}
%[label=a]
{Probabilidad de eventos complementarios}{}
 \centering \Large
 Si la probabilidad de que ocurra un evento es del 40\% la probabilidad de
 que NO ocurra es del 60\%.\\

 \[ P(\text{evento NO ocurra})= 1 - P(\text{evento ocurra}) \]

% \includegraphics[height=.85\textheight]{}
\end{frame}

\begin{frame}
%[label=a]
{Probabilidad de eventos complementarios}{}
 \centering \large
 Sea $ A $ el evento de que al lanzar un dado el resultado sea mayor a 4,
 entonces\\
 $ P(A) = 2/6 $
 \\
 \includegraphics[height=.7\textheight]{eventos-complementarios.png}
\end{frame}

\begin{frame}
%[label=a]
{}{}
 \centering \Large
 ¿Cuál es la probabilidad de que un dado muestre un número par?\\~\\
 \textbf{Ejercicio:} Hacer simulación computacional
% \includegraphics[height=.85\textheight]{}
\end{frame}

\begin{frame}
%[label=a]
{Eventos disjuntos o Mutuamente excluyentes}{}
 \centering
 \begin{itemize}
 \item Eventos disjuntos son sucesos que no pueden ocurrir al mismo tiempo.
 \item No tienen resultados o elementos en común.
 \item Su intersección es igual a cero: $ P(A \cap B) = 0 $
 \end{itemize}

 \vspace{1em}
 \pause
 Ejemplos de eventos disjuntos:

 \begin{itemize}
 \item Lanzar una moneda y obtener cara y sello al mismo tiempo no es
 posible.
 \item Al extraer una sola carta de una baraja no es posible sacar un As y
 una Reina al mismo tiempo.
 \item Obtener un número par y un impar al lanzar un dado una sola vez.
 \end{itemize}
% \includegraphics[height=.85\textheight]{}
\end{frame}

\begin{frame}
%[label=a]
{Probabilidad de eventos disjuntos}{}
 \centering \Large
 Si $ A $ y $ B $ son disjuntos, $ P(A \cup B) = P(A) + P(B) $
% \includegraphics[height=.85\textheight]{}
\end{frame}

\begin{frame}
%[label=a]
{Ejemplo}{}
 \centering \large
% \includegraphics[height=.85\textheight]{}

 ¿Cuál es la probabilidad de que por lo menos, salga una cara en dos tiradas
 de una moneda? \\~\\

 $ E = \{CC , CS, SC, SS\} $\\~\\

 Si asumimos que cada evento ocurre bajo las mismas condiciones, la
 probabilidad viene dada por: \\~\\

 $P(E) = \dfrac{ \{CC, CS, SC\} }{ \{CC , CS, SC, SS\} } = \dfrac{3}{4}$
 \\~\\

 Ejercicio: Corroborar con un experimento computacional
\end{frame}

\begin{frame}
%[label=a]
{Probabilidad de eventos NO disjuntos}{}
 \centering

 Dos eventos son no disjuntos cuando tienen al menos un resultado o elemento
 en común. Si $ A $ y $ B $ son NO disjuntos, entonces

 \[ P(A \cup B) = P(A) + P(B) - P(A \cap B) \]

% \includegraphics[height=.85\textheight]{}
\end{frame}

\begin{frame}
%[label=a]
{Probabilidad de eventos NO disjuntos}{}
 \centering
 \textbf{Ejemplo}\\~\\

 Al lanzar un dado, se quiere calcular la probabilidad de obtener un número
 par (evento $ A $) o un número menor o igual a 2 (evento $ B $).\\~\\

 \pause

 \begin{enumerate}[<+->]

 \item Espacio muestral $E: \{1, 2, 3, 4, 5, 6\}\) (6 resultados
 totales).

 \item Evento $ A $ (número par): \(\{2, 4, 6\} \rightarrow P(A) =
 \frac{3}{6}\).

 \item Evento $ B $ (menor o igual a 2): \(\{1, 2\} \rightarrow P(B) =
 \frac{2}{6}\).

 \item Intersección (\(A \cap B\)): El número 2 está en ambos grupos
 \(\rightarrow P(A \cap B) = \frac{1}{6}\)
 \end{enumerate}

% \includegraphics[height=.85\textheight]{}
\end{frame}

\begin{frame}
%[label=a]
{Independencia}{}
 \centering \Large

 Dos eventos A y B son independientes si la probabilidad de que suceda el
 evento A no cambia la probabilidad del evento B. Estos eventos sí podrían
 suceder al mismo tiempo.
% \includegraphics[height=.85\textheight]{}
\end{frame}

\begin{frame}
%[label=a]
{Probabilidad de eventos independientes}{}
 \centering \large
 Si A y B son eventos independientes, \\~\\

 $ P(A \text{ y } B) = P(A) \cdot P(B) $ \\~\\

 En general, para cualquier número de eventos independientes, la
 probabilidad de que todos los eventos sucedan es el producto de las
 probabilidades de que sucedan los eventos individuales.
\end{frame}

\begin{frame}
%[label=a]
{Ejemplo}{}
 \centering \large
 La probabilidad de que una pareja tenga un hijo y sea mujer es igual a la
 probabilidad de que sea hombre ($ 1/2 $). La probabilidad de que esta
 pareja tenga un hijo con albinismo es $ 1/4 $.\\~\\
 ¿Cuál es la probabilidad de que tengan un hijo varón con albinismo?
% \includegraphics[height=.85\textheight]{}
\end{frame}

\begin{frame}
%[label=a]
{Probabilidad condicional}{}
 \centering \Large
 La probabilidad condicional es la probabilidad de que ocurra un evento
 (llamado $ A $) dado que ya ha sucedido otro evento relacionado ($ B $)\\~\\
  Se calcula mediante la siguiente ecuación:
  \[
  P(A|B)=\frac{P(A \cap B)}{P(B)}
  \]

% \includegraphics[height=.85\textheight]{}
\end{frame}

\begin{frame}
%[label=a]
{Probabilidad condicional}{}
 \centering
 \includegraphics[height=.85\textheight]{Probabilidad_condicionada.pdf}
\end{frame}

\begin{frame}
%[label=a]
{Probabilidad condicional}{}
 \centering
 \includegraphics[height=.85\textheight]{pbb-cond2.jpeg}
\end{frame}

\begin{frame}
%[label=a]
{Ejemplo}{}
 \centering
 Supongamos que en una población de 1000 personas:

 \begin{itemize}
 \item 100 personas tienen una enfermedad (E).

 \item De esas 100 personas con la enfermedad, 80 presentan un síntoma
 específico (S).
 \end{itemize}

 \vspace{1em}
 Queremos saber: ¿cuál es la probabilidad de que una persona tenga la
 enfermedad si sabemos que presenta el síntoma? Es decir: la probabilidad de
 tener la enfermedad dado que se observa el síntoma.\\~\\

 \pause

 Necesitamos conocer:
 \begin{itemize}
 \item Personas con síntoma: supongamos que 200 personas tienen el síntoma.

 \item Personas con enfermedad y síntoma: 80 personas.
 \end{itemize}
% \includegraphics[height=.85\textheight]{}
\end{frame}

\begin{frame}
%[label=a]
{Ejemplo}{}
 \centering
 % Probabilidad condicional
 \[
 P(\text{Enfermedad}|\text{Síntoma})
 \]

 % Aplicación al ejemplo biológico
 \[
 P(E|S)=\frac{P(E \cap S)}{P(S)}
 \]

 % Sustitución con los valores del ejemplo
 \[
 P(E|S)=\frac{80}{200}=0.4
 \]

 En porcentaje:
 \[
 0.4 \times 100 = 40\%
 \]

% \includegraphics[height=.85\textheight]{}
\end{frame}

\begin{frame}
%[label=a]
{}{}
 \centering
 \[
 P(\text{Enfermedad}|\text{Síntoma})=40\%
 \]

 \textbf{Interpretación:} Si una persona presenta el síntoma, la
 probabilidad de que tenga la enfermedad es del 40\%.
% \includegraphics[height=.85\textheight]{}
\end{frame}

\begin{frame}
%[label=a]
{Ejercicio}{}
 \centering
 En un estudio sobre una población de plantas, los investigadores analizan
 la presencia de un gen de resistencia a una enfermedad.\\~\\

 Se estudian 500 plantas y se obtiene la siguiente información:

 \begin{itemize}
 \item 150 plantas poseen el gen de resistencia (R).

 \item De las plantas que poseen el gen de resistencia, 120 sobrevivieron
 después de ser expuestas a la enfermedad (S).

 \item En total, 200 plantas sobrevivieron después de la exposición.
 \end{itemize}

 \vspace{1em}
 ¿Cuál es la probabilidad de que una planta tenga el gen de resistencia
 sabiendo que sobrevivió a la enfermedad?
 \[ P(R|S) \]
% \includegraphics[height=.85\textheight]{}
\end{frame}

%\begin{frame}
%[label=a]
%{}{}
% \centering
% Supongamos que tenemos una bolsa con tres bolas: roja, verde y azul, y
% realizamos el siguiente experimento: Mezclas las bolas y sacas una.
% Después, vuelves a mezclar las bolas y sacas la segunda, sin reponer la
% primera.\\~\\
%
%\begin{itemize}
%\item ¿Cual es la probabilidad de sacar una bola verde primero seguida de
%una roja?
%
%\item ¿Cual es la probabilidad de que uno de las dos bolas sea verde y la
%otra roja?
%
%\end{itemize}
% Realiza el cálculo de forma teórica y comprueba tu solución con un
% experimento computacional.
%
%% \includegraphics[height=.85\textheight]{}
%\end{frame}

\section{Variables aleatorias y distribuciones}

\begin{frame}
%[label=a]
{Variable aleatoria}{}
 \centering \large

 \begin{block}{Variable aleatoria}
 \begin{enumerate}[<+->]
 \item Una variable aleatoria es una forma de representar matemáticamente
 una cantidad cuyo valor no conocemos antes de observarla.

 \item Una variable aleatoria es una variable cuyo valor depende del
 resultado de un proceso aleatorio.
 \end{enumerate}
 \end{block}

% \includegraphics[height=.85\textheight]{}
\end{frame}

\begin{frame}
%[label=a]
{}{}
 \centering \large
 Existen procesos \textbf{fundamentalmente aleatorios} como el decaimiento
 radiactivo\\
 \includegraphics[height=.85\textheight]{decaimiento-radiactivo.jpeg}
\end{frame}

\begin{frame}
%[label=a]
{}{}
 \centering
 En la mayoría de procesos la ``aleatoriedad'' surge del
 \textbf{conocimiento u observabilidad incompleta} del fenómeno.\\
 \vspace{1em}
 \includegraphics[height=.65\textheight]{moneda.jpg}
\end{frame}

\begin{frame}
%[label=a]
{}{}
 \centering
 A partir de una distribución de frecuencias o \textbf{Histograma} podemos
 estimar probabilidades, siguiendo la interpretación frecuentista de la
 probabilidad.\\
 \vspace{.2em}
 \includegraphics[height=.85\textheight]{distribución-de-frecuencias}
\end{frame}

\begin{frame}
%[label=a]
{Distribución de probabilidad}{}
 \centering
 Indica la probabilidad de obtener cierto valor de una \textbf{variable
 aleatoria}
 \begin{columns}
  \column{.5\linewidth}
   \includegraphics[height=.8\textheight]{distribución-pbb}
  \column{.5\linewidth}
  \vspace{-2em}
  \[
  \scalebox{1.5}{$ \displaystyle
  f(x)=\frac{1}{\sigma\sqrt{2\pi}}
  e^{\left(-\frac{1}{2}\left(\frac{x-\mu}{\sigma}\right)^2\right)}
  $}
  \]

%   \includegraphics[width=\linewidth]{}
 \end{columns}

\end{frame}

\begin{frame}
%[label=a]
{}{}
 \centering
 \includegraphics[height=\textheight]{formula-normal-comic.png}
\end{frame}

\begin{frame}
%[label=a]
{}{}
 \centering
 \includegraphics[width=\linewidth]{familia_normales.jpg}
\end{frame}

%\appendix
%\AtNextBibliography{\footnotesize}
%\begin{frame}[plain, allowframebreaks]{Bibliografía}{}
% \bibitemsep 1em
% \printbibliography[heading=none]
%\end{frame}
\end{document}
